\chapter{How to Approach a Problem}\label{ch:approach}

Knowing algorithms is half of competitive programming. The other half is getting from a problem statement to the right
algorithm, under time pressure, without fooling yourself. That second half is a skill like any other: it has a method, the
method can be practised, and people who are good at it are good because they practised it deliberately, not because they
were born knowing which problem wants a segment tree.

This chapter describes that method. It follows one problem from the first reading of its statement to accepted code, then
collects the questions that experienced competitors ask themselves when they are stuck, and ends with the most useful proof
technique for greedy algorithms: the exchange argument.

\begin{chapterintro}{\bronze\ \silver}{nothing; this is the place to start.}
  \item how to read a statement so that it tells you what the solution looks like;
  \item an eight-step loop that takes you from the statement to tested code;
  \item a checklist of questions for when you have no idea;
  \item how to prove a greedy algorithm correct, and how to find the counterexample when it is wrong;
  \item how to tell what kind of problem you are facing, and how to practise.
\end{chapterintro}

% ======================================================================================================================
\section[Anatomy of a problem]{Anatomy of a problem \hfill \bronze}\label{sec:anatomy}

A problem statement is not only a question; it is also a collection of hints about the answer. Every part carries
information, and learning to extract that information is the first skill to build. Figure~\ref{fig:anatomy} dissects a
typical statement.

\begin{figure}[htbp]
\centering
\begin{tikzpicture}
  % one row per part of the statement: the part on the left, what it tells you on the right (rows cannot collide)
  \matrix (m) [matrix of nodes, column sep=10mm, row sep=2.2mm, nodes={anchor=north west, inner sep=3pt},
      column 1/.style={nodes={text width=7.6cm, font=\small, align=left}},
      column 2/.style={nodes={text width=5.4cm, font=\sffamily\scriptsize, text=fgink, align=left}}] {
    |[font=\sffamily\bfseries\normalsize]| Barn Paths & \\
    Farmer John's $n$ barns are joined by $m$ two-way paths. Bessie is in barn 1 and wants to reach barn $n$ walking along
      as few paths as possible. Print that number, or $-1$ if she cannot. &
    \textbf{The story} hides the mathematics: barns are vertices, paths are edges, ``as few paths as possible'' is an
      unweighted shortest path. \\
    \textbf{Input.} The first line has $n$ and $m$. Each of the next $m$ lines has two barns $a_i$ and $b_i$ joined by a
      path. &
    \textbf{The input format} names the objects: an edge list, so a general graph, not a tree and not a grid. \\
    \textbf{Limits.} $2 \le n \le 10^5$, $\ 1 \le m \le 2 \cdot 10^5$. &
    \textbf{The limits} set the budget: about $3 \cdot 10^5$ items, so $\Oh{n + m}$ time. BFS fits. \\
    \textbf{Sample.} \texttt{4 4 / 1 2 / 2 4 / 1 3 / 3 4} $\;\to\;$ \texttt{2} &
    \textbf{The sample} checks your reading: $1 \to 2 \to 4$ uses two paths. \\
    \textbf{Subtask 1 (30 points).} $m = n - 1$ and path $i$ joins barns $i$ and $i + 1$. &
    \textbf{The subtask} is a hint: on a line the answer is $n - 1$. \\
  };
  \begin{scope}[on background layer]
    \node[draw=fgline, fill=fgpaper, rounded corners=3pt, line width=0.6pt, fit=(m-1-1)(m-6-1), inner sep=2pt] {};
  \end{scope}
  \foreach \r in {2,...,6}
    \draw[brace] ($(m-\r-1.north east)+(0.3,-0.05)$) -- ($(m-\r-1.south east)+(0.3,0.05)$);
\end{tikzpicture}
\caption{The parts of a statement and what each one tells you. The limits are the most informative line of all: they
tell you how fast the solution must be before you know what it is.}
\label{fig:anatomy}
\end{figure}

\paragraph{The story.} Almost every statement wraps a mathematical question in a story: cows, cities, robots, games. The
first job is to strip the story away and restate the question in the vocabulary of this book. ``Barns connected by paths''
is a graph. ``A schedule of meetings'' is a set of intervals. ``Pouring water between buckets'' is a state space you can
search. ``Each day the price changes'' is an array, and ``answer $q$ questions about days $l$ to $r$'' is a range query.
Once you can name the objects, you can ask what you know about them.

\paragraph{The input format.} The shape of the input narrows things further. An input with exactly $n - 1$ edges on $n$
vertices is a tree (if it is connected), and trees have their own toolbox (Part~IV). A grid of characters is a graph whose
edges you never store. A list of queries after the data means you must answer many questions fast, which suggests
precomputation or a data structure. Read the input format with the question ``what objects am I given?'' in mind.

\paragraph{The limits.} This is the line that experienced competitors read first. A judge performs about $10^8$ simple
operations per second, so the limits tell you which complexities are possible (Table~\ref{tab:constraints} in the next
chapter). $n \le 20$ invites trying all $2^n$ subsets; $n \le 5000$ allows $\Oh{n^2}$; $n \le 2 \cdot 10^5$ demands
$\Oh{n \log n}$; $n \le 10^{18}$ means the answer is a formula or a logarithmic algorithm. The same story with different
limits is a different problem.

\paragraph{The samples.} Samples exist to check that you understood the question, not to test your code. Work at least one
sample by hand before designing anything. Many wrong answers come from solving a slightly different problem than the one
asked: 0-indexed versus 1-indexed, ``at most'' versus ``exactly'', ``subsequence'' versus ``substring''.

\paragraph{Subtasks.} Olympiads split the points into subtasks with extra restrictions, and those restrictions are hints.
``All values are distinct'', ``the graph is a line'', ``$k = 1$'': each one removes a difficulty, and the solution to a
subtask is often a stepping stone to the full solution. Even when you cannot solve the whole problem, a subtask gives you
points; in the IOI, most medals are decided by partial scores.

\begin{tip}
Before writing code, write the restated problem in one sentence of mathematics at the top of your file as a comment:
\cc{// shortest path from 0 to n-1 in an unweighted graph, n <= 1e5}. If you cannot write that sentence, you have not
understood the problem yet, and no amount of coding will fix that.
\end{tip}

% ======================================================================================================================
\section[A method that works]{A method that works \hfill \bronze}\label{sec:method}

Figure~\ref{fig:method} shows the loop that takes a problem from statement to accepted code. The steps are not new; what
matters is doing all of them, in order, every time. Most people who get stuck on a problem skipped steps 3 and 4 and are
trying to jump straight to a clever idea.

\begin{figure}[htbp]
\centering
\begin{tikzpicture}[stp/.style={fstep=3.4cm, minimum height=1.45cm}]
  % a snake through a 3 x 3 grid: row 1 left to right, row 2 right to left, row 3 left to right
  \node[stp] (s1) at (0, 0) {\textbf{1. Read twice}\\{\footnotesize restate it in one sentence}};
  \node[stp] (s2) at (4.6, 0) {\textbf{2. Limits}\\{\footnotesize which complexity fits?}};
  \node[stp] (s3) at (9.2, 0) {\textbf{3. Small cases}\\{\footnotesize solve by hand, find patterns}};
  \node[stp] (s4) at (9.2, -2.5) {\textbf{4. Brute force}\\{\footnotesize the obviously correct slow solution}};
  \node[stp] (s5) at (4.6, -2.5) {\textbf{5. Find structure}\\{\footnotesize what does the brute force waste?}};
  \node[fq=3.4cm, minimum height=1.45cm] (q) at (0, -2.5) {\textbf{Fast enough?}\\{\footnotesize\mdseries for the limits of step 2}};
  \node[stp] (s6) at (0, -5.0) {\textbf{6. Verify}\\{\footnotesize prove it, or stress test it}};
  \node[stp] (s7) at (4.6, -5.0) {\textbf{7. Code}\\{\footnotesize template, library, clean names}};
  \node[fend=3.4cm, minimum height=1.45cm] (s8) at (9.2, -5.0) {8. Test\\{\footnotesize\mdseries samples, edge cases, largest input}};
  \draw[farr] (s1) -- (s2); \draw[farr] (s2) -- (s3); \draw[farr] (s3) -- (s4);
  \draw[farr] (s4) -- (s5); \draw[farr] (s5) -- (q);
  \draw[farr] (q) -- node[fyes, right] {yes} (s6);
  \draw[farr] (s6) -- (s7); \draw[farr] (s7) -- (s8);
  % the two ways back to step 5, each in its own gap between the rows
  \draw[farr, fgverm] ($(q.north)+(0.9,0)$) to[out=65, in=115] node[fno, above] {no} ($(s5.north)+(-0.9,0)$);
  \draw[farr, fgverm, dashed] ($(s6.north)+(1.2,0)$) to[out=60, in=-120] node[fno, pos=0.55, right=2pt] {counterexample}
    ($(s5.south)+(-1.0,0)$);
\end{tikzpicture}
\caption{The problem-solving loop. Steps 3 to 5 are where the thinking happens; step 6 is what separates people who get
accepted from people who get \emph{wrong answer on test 17}.}
\label{fig:method}
\end{figure}

We follow one problem through all eight steps. It is a classic from the CSES Problem Set, and it is a good example
because the right idea is not visible in the statement.

\begin{problem}{Array Division (CSES)}
You are given $n$ positive integers and a number $k$. Divide the array into $k$ contiguous, non-empty parts so that the
largest sum of a part is as small as possible, and print that largest sum. Limits: $1 \le k \le n \le 2 \cdot 10^5$,
values up to $10^9$.
\end{problem}

\subsection{Step 1: read twice, restate once}

Read the whole statement, then restate it in one sentence of mathematics: \emph{choose $k - 1$ cut positions in the array
to minimize the maximum part sum.} Check the restatement against the words: the parts are contiguous (so we cut, we do not
choose subsets), every part is non-empty (so the cuts are distinct), and all values are positive (remember this: positivity
often matters). The sample in CSES is $n = 5$, $k = 3$, array $2, 4, 7, 3, 5$, answer 8.

\subsection{Step 2: let the limits talk}

$n \le 2 \cdot 10^5$ rules out anything quadratic: $n^2 = 4 \cdot 10^{10}$ operations would take minutes. We need roughly
$\Oh{n \log n}$ or $\Oh{n \log(\text{something})}$. That already hints at sorting, a binary search, a heap or a segment
tree, and it rules out ``try every way to cut'', which is $\binom{n-1}{k-1}$, astronomically many. The sums can reach
$2 \cdot 10^5 \cdot 10^9 = 2 \cdot 10^{14}$, so the answer needs 64-bit integers.

\subsection{Step 3: solve small cases by hand}

Take the sample and try splits by hand (Figure~\ref{fig:arraydiv-cases}). A few observations fall out immediately. The
answer is at least the largest element (7 here), because that element sits in some part. It is at most the total (21),
the answer for $k = 1$. And the good splits are ``balanced'': moving a cut towards the heavier side makes things better.

\begin{figure}[htbp]
\centering
\begin{tikzpicture}[psum/.style={font=\sffamily\scriptsize, text=fggray}, pmax/.style={font=\sffamily\scriptsize\bfseries, text=fgverm}]
  % each row: the array, the two cuts (bars), the sum of every part below it, the largest sum on the right
  \foreach \y [count=\r] in {0, -1.55, -3.1, -4.65} \arr{R\r}{0}{\y}{2,4,7,3,5};
  \newcommand{\arrcut}[2]{\draw[line width=2.4pt, fgverm] ({#1*0.7}, {#2+0.5}) -- ({#1*0.7}, {#2-0.5});}
  \arrcut{2}{0} \arrcut{3}{0} \arrcut{3}{-1.55} \arrcut{4}{-1.55} \arrcut{1}{-3.1} \arrcut{3}{-3.1} \arrcut{2}{-4.65} \arrcut{4}{-4.65}
  % part sums: centred under each part
  \foreach \x/\y/\s/\st in {0.7/0/6/psum, 1.75/0/7/psum, 2.8/0/8/pmax,
      1.05/-1.55/13/pmax, 2.45/-1.55/3/psum, 3.15/-1.55/5/psum,
      0.35/-3.1/2/psum, 1.4/-3.1/11/pmax, 2.8/-3.1/8/psum,
      0.7/-4.65/6/psum, 2.1/-4.65/10/pmax, 3.15/-4.65/5/psum}
    \node[\st] at (\x, {\y-0.58}) {\s};
  \foreach \y/\m/\c in {0/8/fggreen, -1.55/13/fggray, -3.1/11/fggray, -4.65/10/fggray}
    \node[font=\sffamily\small, text=\c, anchor=west] at (4.1, \y) {largest part: \textbf{\m}};
  \node[font=\sffamily\scriptsize\bfseries, text=fggreen, anchor=east] at (-0.25, 0) {best};
  \begin{scope}[on background layer]
    \fill[fggreen!9, rounded corners=3pt] (-1.0, 0.62) rectangle (6.9, -0.82);
  \end{scope}
\end{tikzpicture}
\caption{Four of the six ways to cut $2, 4, 7, 3, 5$ into three parts. The best one (green) has maximum 8; no split can do
better than 7, the largest element, and here even 7 is impossible.}
\label{fig:arraydiv-cases}
\end{figure}

\subsection{Step 4: write down the brute force}

The obviously correct solution tries every way to cut: choose $k - 1$ of the $n - 1$ gaps and take the best. A slightly
smarter brute force is a DP, $\mathrm{best}[i][j]$ = the smallest possible maximum when the first $i$ elements are cut into
$j$ parts, which takes $\Oh{n^2 k}$ time. Both are far too slow for the real limits, and that is fine: the brute force is
not meant to be submitted. It fixes \emph{what} we compute, gives us a reference to test against, and for small subtasks it
is often worth points on its own.

\subsection{Step 5: find structure}

Now ask what the brute force wastes. It builds every split and \emph{then} measures its maximum. What if we turn the
question around and fix the maximum first? ``Can the array be cut into $k$ parts, each with sum at most $C$?'' is a much
easier question, because a greedy answers it: walk through the array, extend the current part while it still fits under
$C$, and start a new part when it does not. Extending a part as far as possible never hurts, since it leaves the fewest
elements for the remaining parts (this is an exchange argument, Section~\ref{sec:exchange}). If the greedy needs at most $k$
parts, the answer is yes; with fewer than $k$ parts we can always split some part further, because every value is positive
and every part has at least one element.

The question also has the crucial property of \emph{monotonicity}: if cap $C$ works, every larger cap works too. So the
answers form a block of ``no'' followed by a block of ``yes'' (Figure~\ref{fig:arraydiv-mono}), and we can binary search for
the boundary. This is binary search on the answer, the subject of Chapter~\ref{ch:binsearch}.

\begin{figure}[htbp]
\centering
\begin{tikzpicture}
  \arr{C}{0}{0}{7,8,9,10,11,12,13}
  \node[lbl, anchor=east] at (-0.15, 0) {cap $C$};
  \arr{P}{0}{-0.85}{4,3,3,3,3,3,2}
  \node[lbl, anchor=east] at (-0.15, -0.85) {parts needed};
  \arrs{V}{0}{-1.7}{no/hot, yes/good, yes/good, yes/good, yes/good, yes/good, yes/good}
  \node[lbl, anchor=east] at (-0.15, -1.7) {$\le k = 3$?};
  \draw[hded] ($(V-2.south)+(0,-0.7)$) node[below, lbl, text=fgverm] {smallest cap that works: 8} -- (V-2.south);
  \node[notel, text width=6.3cm, anchor=west] at (5.4, -0.85)
    {Greedy with cap 8 cuts $2,4 \mid 7 \mid 3,5$. With cap 7 it needs $2,4 \mid 7 \mid 3 \mid 5$: four parts, too many.
     Larger caps never need more parts, so the column of answers is ``no'' then ``yes''.};
\end{tikzpicture}
\caption{For each cap $C$, the number of parts the greedy needs. The yes/no row is monotone, so binary search finds the
first ``yes'' in about $\log_2(2\cdot 10^{14}) \approx 48$ steps.}
\label{fig:arraydiv-mono}
\end{figure}

\subsection{Step 6: verify before you code}

Two claims carry the solution: the greedy check is correct, and the answers are monotone. The monotonicity is immediate (a
split that respects cap $C$ also respects $C + 1$). For the greedy, a short exchange argument works: take any valid split
into at most $k$ parts; its first part ends no later than the greedy's first part, because the greedy extends as far as the
cap allows; moving that first cut to where the greedy put it only shrinks the second part; repeat. So if any valid split
exists, the greedy finds one with no more parts. When you cannot find a proof in a few minutes, the fallback is to stress
test the idea against the brute force on thousands of small random inputs (Section~\ref{sec:stress-idea}).

\subsection{Step 7: code}

Only now do we type. The search range is from the largest element (the answer can never be smaller) to the total (the
answer for one part).

\begin{code}
void solve() {
    int n, k;
    cin >> n >> k;
    vi a(n);
    cin >> a;
    auto parts_needed = [&](ll cap) {  // greedy: extend a part while it fits
        ll parts = 1, cur = 0;
        for (ll x : a) {
            if (cur + x > cap) parts++, cur = 0;
            cur += x;
        }
        return parts;
    };
    ll lo = *max_element(all(a)), hi = accumulate(all(a), 0LL);
    while (lo < hi) {  // smallest cap that needs at most k parts
        ll mid = lo + (hi - lo) / 2;
        if (parts_needed(mid) <= k) hi = mid;
        else lo = mid + 1;
    }
    print(lo);
}
\end{code}

The complexity is $\Oh{n \log S}$ where $S$ is the total: about $2 \cdot 10^5 \cdot 48 \approx 10^7$ operations, which runs
in a few milliseconds. This code, like every complete program in this book, was compiled with your template and checked
against a brute force on random inputs.

\subsection{Step 8: test}

Run the CPH samples first. Then think about the edges of the input space: $n = 1$; $k = 1$ (the answer is the total);
$k = n$ (the answer is the maximum); all elements equal; one huge element among small ones; the largest possible input (to
check speed and 64-bit overflow). Each of these catches a different class of bug. If anything is uncertain, a stress test
against the brute force settles it.

\begin{fact}[The loop in real contests]
Experienced competitors run through this loop in their heads in a minute or two for easy problems, and spend most of their
contest time in steps 3 to 6 for hard ones. The two most common causes of a lost problem are skipping step 2 (an $\Oh{n^2}$
solution for $n = 2 \cdot 10^5$) and skipping step 6 (a greedy that ``obviously'' works and does not).
\end{fact}

% ======================================================================================================================
\section[Thinking tools]{Thinking tools \hfill \bronze}\label{sec:thinking}

Step 5 of the loop, ``find structure'', is where problems are solved, and it is the step that feels least like a method.
It becomes much more mechanical once you have a set of moves to try. These are the ones that work most often, adapted from
George P\'olya's classic \emph{How to Solve It}.

\paragraph{Draw a picture.} Arrays become bars, intervals become segments on a line, pairs $(x, y)$ become points in the
plane, states become vertices. A picture makes monotonicity, overlaps and symmetry visible in a way that indices never do.
Most figures in this book exist because the author needed them to see the idea.

\paragraph{Specialise.} Solve an easier version first: all values equal, $k = 1$, the graph is a path, the grid has one row.
The special case either solves the problem outright (surprisingly often) or shows which part is hard. Subtasks are special
cases chosen by the problem setter, and they are chosen because they are stepping stones.

\paragraph{Generalise.} Sometimes the more general problem is easier. ``Is there a path from 1 to $n$?'' is best answered by
computing \emph{all} reachable vertices. ``Answer this for $k = 5$'' becomes easier as ``answer this for every $k$'', which
may have a recurrence.

\paragraph{Work backwards.} Start from the goal. What must the last move of a game look like? What does the final state of a
process look like, and what could have come just before it? Many DP and game problems are clearest when read from the end.

\paragraph{Look at the answer.} What does an optimal solution look like? Can you assume it has a normal form (sorted, with no
crossing pairs, with the largest element first)? Proving that ``some optimal solution has form X'' often cuts the search
space from exponential to polynomial.

\paragraph{Change one thing at a time.} If the answer for an input is known, how does it change when one element is added,
removed or modified? The answer to that question is often an incremental algorithm or a data structure.

\paragraph{Find the bottleneck.} In the brute force, which part is slow? If it is ``for every pair, compute something'', can
the something be precomputed? If it is ``for every query, scan the array'', can the array be preprocessed? Improving the
bottleneck by one factor of $n$ is usually the whole solution.

% ======================================================================================================================
\section[When you are stuck]{When you are stuck: a checklist \hfill \silver}\label{sec:stuck}

When nothing comes, go through these questions one by one. Each has solved thousands of problems, and together they cover
the large majority of what appears in contests up to the Gold level. Chapter~\ref{ch:patterns} develops the most important
of them into full patterns with worked examples.

\begin{description}[style=nextline, leftmargin=1.2em, itemsep=4pt, font=\sffamily\bfseries\small\color{cpblue}]
  \item[Is the answer monotone in some parameter?] If ``$x$ works'' implies ``$x + 1$ works'', binary search on the answer
    turns an optimisation problem into a sequence of yes/no checks, as in Array Division above (Chapter~\ref{ch:binsearch}).
  \item[Would sorting change nothing?] If the answer does not depend on the input order, sort it. Sorted data enables greedy
    choices, two pointers, sweeps and binary search (Chapter~\ref{ch:sorting}).
  \item[Can each element's contribution be counted separately?] ``Sum over all subarrays of their minimum'' is hard as
    stated; ``for each element, in how many subarrays is it the minimum?'' is easy with a monotonic stack. This is the
    contribution technique (Section~\ref{sec:contribution}).
  \item[Is the complement easier?] Counting ``at least one'' is often done as total minus ``none''. Counting pairs that are
    \emph{not} coprime is total minus coprime pairs. Inclusion--exclusion generalises this (Chapter~\ref{ch:combinatorics}).
  \item[Can I fix one thing and solve the rest fast?] Fix the left endpoint, the maximum element, the middle of a
    palindrome, the first move. Then the remaining problem is often one-dimensional and fast with a data structure.
  \item[Can I process things in a different order?] Offline (read all queries first, answer them in a convenient order),
    backwards in time (deletions become insertions, which DSU handles), or sorted by some key (sweep line).
  \item[Is there an invariant?] A quantity that never changes proves that something is impossible; a quantity that always
    decreases proves that a process terminates. Parity and colouring arguments are the classic ones
    (Section~\ref{sec:invariants}).
  \item[Does the answer depend only on a small state?] If the future depends only on ``where I am'' and a few numbers, not
    on the whole history, dynamic programming applies (Part~V).
  \item[Are things connected to other things?] Model it as a graph: vertices for the objects, edges for the relation, and
    reach for Part~III.
  \item[Is some parameter small?] Values up to 100, $k \le 20$, at most 60 distinct colours: iterate over the small
    parameter, or use it as a bitmask.
  \item[Do large and small cases behave differently?] Handle ``heavy'' items (few, each large) and ``light'' items (many,
    each small) with different methods; the threshold $\sqrt n$ balances them (Chapter~\ref{ch:static}).
  \item[What does an optimal solution look like?] Prove it can be assumed sorted, non-crossing or canonical, then search only
    among such solutions.
  \item[Is it a known problem in disguise?] Shortest path, matching, interval scheduling, LIS, knapsack and flow appear in
    costume constantly. The recognition table in Chapter~\ref{ch:craft} maps statement signals to techniques.
\end{description}

\begin{trick}[Change the question]
Many hard problems become easy when you ask a different but equivalent question. ``Minimise the maximum'' becomes ``can the
maximum be at most $X$?''. ``Count pairs with property $P$'' becomes ``for each element, count the partners seen so far''.
``Is there a path?'' becomes ``are they in the same component?''. ``What is the $k$-th smallest?'' becomes ``how many are at
most $x$?''. Train yourself to rephrase every problem at least twice before choosing an algorithm; the second phrasing is
usually the one with a known solution.
\end{trick}

% ======================================================================================================================
\section[Proving greedy choices]{Proving greedy choices: the exchange argument \hfill \silver}\label{sec:exchange}

A greedy algorithm builds a solution one piece at a time, always taking the choice that looks best right now, and never
reconsidering. Greedy algorithms are short, fast and very often wrong. The exchange argument is the standard way to tell
the difference before you spend half an hour coding one.

\subsection{The argument}

To show that a greedy algorithm is optimal, show that \emph{some} optimal solution agrees with the greedy's first choice.
Then the rest of the problem is a smaller instance of the same kind, and induction finishes the proof:

\begin{enumerate}
  \item Take any optimal solution $O$. If its first choice equals the greedy choice $g$, there is nothing to do.
  \item Otherwise, \emph{exchange}: modify $O$ so that it makes the choice $g$, and show that the result $O'$ is still valid
        and is no worse than $O$. Then $O'$ is optimal too.
  \item $O'$ agrees with the greedy on the first choice; repeat the argument on what remains.
\end{enumerate}

\begin{problem}{Interval scheduling}
Given $n$ intervals, choose as many pairwise disjoint ones as possible.
\solution Sort by end time and repeatedly take the first interval that does not overlap the ones already taken. Exchange:
let $O$ be an optimal solution whose first interval $o_1$ is not the greedy's first interval $g$, which is the interval that
ends first overall. Replace $o_1$ by $g$ (Figure~\ref{fig:exchange-intervals}). Since $g$ ends no later than $o_1$, it
cannot overlap $o_2$, which starts after $o_1$ ends. The new solution is valid and has the same size, so it is optimal and
starts with $g$.
\usetemplate{Intervals: merge + max disjoint}{Ctrl+Alt+T, Shift+N}{cpintervals}
\end{problem}

\begin{figure}[htbp]
\centering
\begin{tikzpicture}[x=0.72cm]
  \numline{0}{12}{0}{time}
  \node[lbl, anchor=east] at (-0.4, 1.9) {optimal $O$};
  \ivl{1}{5}{1.9}{fgline}{$o_1$} \ivl{6}{8}{1.9}{fgblue!55}{$o_2$} \ivl{9}{12}{1.9}{fgblue!55}{$o_3$}
  \node[lbl, anchor=east] at (-0.4, 0.85) {exchanged $O'$};
  \ivl{2}{4}{0.85}{fggreen!70}{$g$} \ivl{6}{8}{0.85}{fgblue!55}{$o_2$} \ivl{9}{12}{0.85}{fgblue!55}{$o_3$}
  \draw[xed] (4, 0.45) -- (4, 2.35); \draw[xed] (5, 0.45) -- (5, 2.35);
  \node[notel, anchor=west, text width=3.7cm] at (13.1, 1.4) {$g$ ends at 4, before $o_1$ ends at 5, so it cannot reach
    $o_2$. Same size, still valid.};
\end{tikzpicture}
\caption{The exchange for interval scheduling: swapping the first interval of an optimal solution for the interval that ends
earliest keeps the solution valid and just as large.}
\label{fig:exchange-intervals}
\end{figure}

\subsection{A second example: adjacent swaps}

When the greedy is an \emph{ordering} (``process items sorted by $X$''), the cleanest exchange is between two neighbours.
Take any order that is not sorted by the greedy key; somewhere two adjacent items are in the wrong order; swapping them
changes only their own contribution, so you just compare two numbers.

\begin{problem}{Tasks and Deadlines (CSES)}
There are $n$ tasks; task $i$ takes $a_i$ time units and has deadline $d_i$. You do them one after another, in an order you
choose, starting at time 0. Finishing task $i$ at time $f_i$ earns $d_i - f_i$ points (possibly negative). Maximise the
total. Limits: $n \le 2 \cdot 10^5$.
\solution The total is $\sum d_i - \sum f_i$, and $\sum d_i$ does not depend on the order, so we must minimise the sum of
finishing times. Do the shortest task first. Proof by adjacent swap (Figure~\ref{fig:exchange-swap}): if task $i$ comes
immediately before task $j$ and both start after time $T$, their finishing times are $T + a_i$ and $T + a_i + a_j$; swapped,
they are $T + a_j$ and $T + a_j + a_i$. Every other task is unaffected, so the swap changes the total by $a_j - a_i$. If
$a_i > a_j$, swapping strictly helps; hence an optimal order has no ``longer before shorter'' neighbours, which means it is
sorted by duration.
\end{problem}

\begin{figure}[htbp]
\centering
\begin{tikzpicture}[x=0.75cm]
  \numline{0}{11}{0}{time}
  \node[lbl, anchor=east] at (-0.5, 1.9) {long first};
  \fill[fgline!40] (0, 1.72) rectangle (3, 2.08); \node[font=\sffamily\scriptsize] at (1.5, 1.9) {earlier};
  \fill[fgverm!45] (3, 1.72) rectangle (8, 2.08); \node[font=\sffamily\scriptsize] at (5.5, 1.9) {$i$: $a_i = 5$};
  \fill[fgblue!40] (8, 1.72) rectangle (10, 2.08); \node[font=\sffamily\scriptsize] at (9, 1.9) {$j$};
  \node[lbl, anchor=west] at (11.2, 1.9) {finish at 8 and 10: sum 18};
  \node[lbl, anchor=east] at (-0.5, 0.85) {short first};
  \fill[fgline!40] (0, 0.67) rectangle (3, 1.03); \node[font=\sffamily\scriptsize] at (1.5, 0.85) {earlier};
  \fill[fgblue!40] (3, 0.67) rectangle (5, 1.03); \node[font=\sffamily\scriptsize] at (4, 0.85) {$j$};
  \fill[fgverm!45] (5, 0.67) rectangle (10, 1.03); \node[font=\sffamily\scriptsize] at (7.5, 0.85) {$i$: $a_i = 5$};
  \node[lbl, anchor=west] at (11.2, 0.85) {finish at 5 and 10: sum 15};
  \draw[xed] (10, 0.5) -- (10, 2.3);
\end{tikzpicture}
\caption{Swapping two adjacent tasks leaves everything outside them untouched (both orders end at time 10), and puts the
shorter task's finishing time first. The sum drops by $a_i - a_j = 3$.}
\label{fig:exchange-swap}
\end{figure}

\begin{code}
void solve() {
    int n;
    cin >> n;
    vector<pll> task(n);  // (duration, deadline)
    for (auto& [a, d] : task) cin >> a >> d;
    sort(all(task));  // shortest first
    ll now = 0, reward = 0;
    for (auto [a, d] : task) now += a, reward += d - now;
    print(reward);
}
\end{code}

Notice what the deadlines did: nothing. The statement made them look important, the mathematics showed they are irrelevant.
This happens more often than you would think, and it is exactly what the restatement in step 1 is for.

\subsection{When greedy is wrong}

A greedy that ``obviously works'' fails more often than it succeeds. The classic example is making change: with coins
$\{1, 3, 4\}$ and amount 6, ``always take the largest coin that fits'' pays $4 + 1 + 1$ (three coins), but $3 + 3$ uses two.
The exchange argument breaks exactly where the greedy breaks: you cannot swap the 4 into an optimal solution without paying
extra elsewhere. The cure is dynamic programming (Chapter~\ref{ch:dp}).

\begin{tip}[Hunting counterexamples]
If you cannot prove a greedy in a few minutes, look for a counterexample in the smallest cases: two or three items, values
from 1 to 5. Wrong greedy ideas almost always fail on inputs of size 3 to 6. Even better, let the computer look: a brute
force plus a random generator plus a loop finds the counterexample in seconds (next section).
\end{tip}

% ======================================================================================================================
\section[Testing an idea before trusting it]{Testing an idea before trusting it \hfill \silver}\label{sec:stress-idea}

A \defn{stress test} runs your fast solution and a brute force on thousands of small random inputs and stops at the first
input where they disagree (Figure~\ref{fig:stress}). It is the single most effective debugging technique in competitive
programming, and it is just as useful \emph{before} coding the real solution: implement the brute force and a quick,
throwaway version of the idea, and let the stress test tell you whether the idea is right.

\begin{figure}[htbp]
\centering
\begin{tikzpicture}
  \node[fstep=2.7cm, minimum height=1.3cm] (gen) at (0, 0) {\textbf{generator}\\{\footnotesize random small input from seed $s$}};
  \node[fstep=2.3cm] (fast) at (4.1, 1.05) {\textbf{fast}\\{\footnotesize your solution}};
  \node[fstep=2.3cm] (brute) at (4.1, -1.05) {\textbf{brute}\\{\footnotesize obviously correct}};
  \node[fq=2.1cm, minimum height=1.1cm] (cmp) at (7.9, 0) {\textbf{same output?}};
  \node[fend=3.1cm, minimum height=1.3cm] (stop) at (11.9, 0) {counterexample\\{\footnotesize\mdseries the smallest failing input}};
  \draw[farr] (gen.east) -- (fast.west); \draw[farr] (gen.east) -- (brute.west);
  \draw[farr] (fast.east) -- (cmp.west); \draw[farr] (brute.east) -- (cmp.west);
  \draw[farr] (cmp) -- node[fno, above] {no} (stop);
  \draw[farr] (cmp.south) -- (7.9, -2.3) -- node[fyes, below] {yes: try the next seed $s + 1$} (0, -2.3) -- (gen.south);
\end{tikzpicture}
\caption{A stress test. Because the inputs are small, the counterexample it finds is small enough to debug by hand. Your
setup runs this loop with \key{Ctrl+Alt+X} (GUIDE.pdf, Section~4.4).}
\label{fig:stress}
\end{figure}

Three details make stress tests effective. First, keep the inputs \emph{tiny}: $n$ from 1 to 6, values from 1 to 10. Bugs
show up on small inputs just as reliably as on large ones, and a failing input of size 4 can be debugged by hand. Second,
make the brute force as dumb as possible, so that it is obviously correct; its speed does not matter. Third, when the test
finds a counterexample, add it to CPH as a test case before fixing anything, so that it stays fixed.

% ======================================================================================================================
\section[Kinds of problems]{Kinds of problems \hfill \bronze}\label{sec:kinds}

Problems fall into a handful of kinds, and each kind has its own habits. Recognising the kind is often the first real step
towards the solution.

\begin{center}
\small
\begin{tabularx}{\linewidth}{p{2.6cm}Lp{4.8cm}}
\toprule
\thd{Kind} & \thd{Asks for} & \thd{Typical tools} \\
\midrule
Decision & yes or no: is it possible? & invariants, greedy, graph search, 2-SAT \\
Optimisation & the best value, or a best object & greedy, DP, binary search on the answer, flows \\
Counting & how many (usually modulo $10^9+7$) & combinatorics, DP, inclusion--exclusion \\
Construction & any valid object, or proof there is none & patterns from small cases, induction, greedy \\
Queries & many questions on one structure & prefix sums, segment trees, offline processing \\
Interactive & ask questions, deduce a hidden object & binary search, adaptive strategies \\
Games & who wins with perfect play & game DP, Sprague--Grundy, symmetry \\
\bottomrule
\end{tabularx}
\end{center}

\paragraph{Decision problems} often have a clean characterisation: ``possible if and only if the sum is even and no element
exceeds half of it''. Look for the invariant that makes impossible cases impossible, then check that everything else is
possible.

\paragraph{Optimisation problems} ask for a minimum or maximum. The first question is always whether a greedy works (try to
prove it, or break it on small cases); the second is whether the answer is monotone in some parameter (binary search); the
third is what state a DP would need.

\paragraph{Counting problems} print huge numbers modulo a prime. They almost never ask you to enumerate: count by
structure. Decompose the objects you count, count each part independently, multiply; when parts overlap, use
inclusion--exclusion; when the structure is sequential, use DP.

\paragraph{Construction problems} accept any valid answer. Solve small cases by hand and look for a pattern that
generalises, often by induction: a construction for $n$ from one for $n - 1$ or $n - 2$. Checking your output with a small
validator program is worth the minute it takes.

\paragraph{Query problems} give data and many questions. Decide first whether the data changes (static or dynamic), whether
the queries can be answered offline (read them all first), and what operation combines two ranges (sum, min, gcd). Part~VI is
about exactly this.

\paragraph{Interactive problems} hide an object and let you ask questions. The number of allowed questions is the key
constraint: $\log_2$ of the number of possibilities means binary search; more questions allow more naive strategies
(Chapter~\ref{ch:special}).

\paragraph{Games} ask who wins with perfect play. Small games are solved by DP over positions (winning if some move leads to
a losing position); sums of independent games are solved by Sprague--Grundy theory (Chapter~\ref{ch:combinatorics}).

% ======================================================================================================================
\section[How to practise]{How to practise \hfill \bronze}\label{sec:practise}

Solving problems is the only way to get better at solving problems, but not all practice is equal.

\begin{itemize}
  \item \textbf{Work slightly above your level.} Problems you solve in five minutes teach little; problems you cannot even
        start teach little too. The sweet spot is problems you solve in 30 to 60 minutes, or solve after reading one hint.
  \item \textbf{Upsolve.} After every contest, solve the problems you could not solve during it, reading the editorial if
        needed. Then write down, in one sentence, the idea you were missing. This habit is worth more than any single
        technique in this book.
  \item \textbf{Implement, do not just read.} Reading an editorial gives the feeling of understanding; implementing it shows
        whether you really do. Most of the details that matter only appear when you write the code.
  \item \textbf{Build a library and know it.} This book comes with one (GUIDE.pdf). Knowing what is in it, and pasting it
        with one shortcut, saves minutes in every contest and removes a whole class of bugs.
  \item \textbf{Review your mistakes.} Keep a list of the bugs you made (overflow, off-by-one, uninitialised variables,
        wrong reading of the statement). The same five bugs cause most of everyone's wrong answers.
\end{itemize}

\begin{pitfall}[Coding before thinking]
Five minutes of thinking saves thirty minutes of debugging. If you cannot explain your algorithm, and why it is correct, in
three sentences, do not start typing. The pressure of a contest makes typing feel like progress; it usually is not.
\end{pitfall}

\begin{summary}[Chapter summary]
\begin{itemize}
  \item Read the statement for information: the story names the objects, the input format their shape, the limits the
        complexity, the samples your understanding, the subtasks the path to the solution.
  \item Follow the loop: read, budget, small cases, brute force, structure, verify, code, test. Do not skip steps 3, 4 and 6.
  \item When stuck, run through the checklist; the most powerful move is to change the question.
  \item Prove greedy algorithms with an exchange argument, or break them with a small counterexample, ideally found by a
        stress test.
\end{itemize}
\end{summary}

\begin{practice}
\ppgroup{Warm-up: reading and restating}
\begin{pplist}
  \item Weird Algorithm
  \item Missing Number
  \item Repetitions
  \item Increasing Array
  \item Permutations
  \item Number Spiral
  \item Two Knights
  \item Two Sets
  \item Coin Piles
\end{pplist}
\ppgroup{Small cases, patterns and constructions}
\begin{pplist}
  \item Bit Strings
  \item Trailing Zeros
  \item Palindrome Reorder
  \item Gray Code
  \item Tower of Hanoi
  \item Digit Queries
\end{pplist}
\ppgroup{Greedy and exchange arguments}
\begin{pplist}
  \item Stick Lengths
  \item Movie Festival
  \item Tasks and Deadlines
  \item Array Division
  \item Ferris Wheel
  \item Missing Coin Sum
\end{pplist}
\end{practice}
